\documentclass{article}
\usepackage{amsmath, amssymb, amsthm}
\usepackage{geometry}
\geometry{a4paper, margin=1in}
\usepackage{fancyhdr}
\usepackage{hyperref}
\usepackage[english, russian]{babel}
\usepackage[utf8]{inputenc}
\usepackage{listings}
\usepackage{xcolor}
\usepackage{tikz}

\pagestyle{fancy}
\fancyhf{}
\fancyhead[L]{Algorithms}
\fancyhead[R]{\thepage}

\title{Algorithms 1 \\ Homework 9}
\author{Nikita Zagainov}
\date{\selectlanguage{english}\today}

\lstset{
    language=Python,
    basicstyle=\ttfamily\small,
    keywordstyle=\color{blue},
    stringstyle=\color{red},
    commentstyle=\color{gray},
    backgroundcolor=\color{lightgray!20},
    frame=single,
    rulecolor=\color{black},
    numbers=left,
    numberstyle=\tiny\color{gray},
    stepnumber=1,
    numbersep=10pt,
    showspaces=false,
    showstringspaces=false,
    breaklines=true,
    breakatwhitespace=true,
    tabsize=4,
    captionpos=b
}

\begin{document}

\maketitle

\begin{center}
	\section*{Кратчайшие пути в графах}
\end{center}
\begin{enumerate}
	\item Шаг 0:
	      \begin{center}
		      \begin{tikzpicture}
			      \node[draw, circle] (e) at (0,2) {e};
			      \node[draw, circle] (b) at (2,2) {b};
			      \node[draw, circle] (a) at (4,2) {a};
			      \node[draw, circle] (d) at (6,2) {d};
			      \node[draw, circle] (h) at (0,0) {h};
			      \node[draw, circle] (g) at (3,0) {g};
			      \node[draw, circle] (f) at (6,0) {f};

			      \node[above] at (e.north) {$0$};
			      \node[above] at (b.north) {$\infty$};
			      \node[above] at (a.north) {$\infty$};
			      \node[above] at (d.north) {$\infty$};
			      \node[above] at (h.north) {$\infty$};
			      \node[below] at (g.south) {$\infty$};
			      \node[below] at (f.south) {$\infty$};

			      \draw[->, bend left=15, >=stealth, line width=1pt] (e) to node[above] {-1} (b);
			      \draw[->, bend left=15, >=stealth, line width=1pt] (b) to node[above] {4} (a);
			      \draw[->, >=stealth, line width=1pt] (e) to node[above] {-2} (g);
			      \draw[->, >=stealth, line width=1pt] (h) to node[below] {-3} (g);
			      \draw[->, >=stealth, line width=1pt] (g) to node[below] {1} (f);
			      \draw[->, bend left=15, >=stealth, line width=1pt] (f) to node[left] {2} (a);
			      \draw[->, bend right=15, >=stealth, line width=1pt] (d) to node[left] {4} (f);
			      \draw[->, bend right=15, >=stealth, line width=1pt] (f) to node[right] {-3} (d);
		      \end{tikzpicture}
	      \end{center}

	      Шаг 1:
	      \begin{center}
		      \begin{tikzpicture}
			      \node[draw, circle] (e) at (0,2) {e};
			      \node[draw, circle] (b) at (2,2) {b};
			      \node[draw, circle] (a) at (4,2) {a};
			      \node[draw, circle] (d) at (6,2) {d};
			      \node[draw, circle] (h) at (0,0) {h};
			      \node[draw, circle] (g) at (3,0) {g};
			      \node[draw, circle] (f) at (6,0) {f};

			      \node[above] at (e.north) {$0$};
			      \node[above] at (b.north) {$-1$};
			      \node[above] at (a.north) {$\infty$};
			      \node[above] at (d.north) {$\infty$};
			      \node[above] at (h.north) {$\infty$};
			      \node[below] at (g.south) {$-2$};
			      \node[below] at (f.south) {$\infty$};

			      \draw[->, bend left=15, >=stealth, line width=1pt] (e) to node[above] {-1} (b);
			      \draw[->, bend left=15, >=stealth, line width=1pt] (b) to node[above] {4} (a);
			      \draw[->, >=stealth, line width=1pt] (e) to node[above] {-2} (g);
			      \draw[->, >=stealth, line width=1pt] (h) to node[below] {-3} (g);
			      \draw[->, >=stealth, line width=1pt] (g) to node[below] {1} (f);
			      \draw[->, bend left=15, >=stealth, line width=1pt] (f) to node[left] {2} (a);
			      \draw[->, bend right=15, >=stealth, line width=1pt] (d) to node[left] {4} (f);
			      \draw[->, bend right=15, >=stealth, line width=1pt] (f) to node[right] {-3} (d);
		      \end{tikzpicture}
	      \end{center}

	      Шаг 2:
	      \begin{center}
		      \begin{tikzpicture}
			      \node[draw, circle] (e) at (0,2) {e};
			      \node[draw, circle] (b) at (2,2) {b};
			      \node[draw, circle] (a) at (4,2) {a};
			      \node[draw, circle] (d) at (6,2) {d};
			      \node[draw, circle] (h) at (0,0) {h};
			      \node[draw, circle] (g) at (3,0) {g};
			      \node[draw, circle] (f) at (6,0) {f};

			      \node[above] at (e.north) {$0$};
			      \node[above] at (b.north) {$-1$};
			      \node[above] at (a.north) {$3$};
			      \node[above] at (d.north) {$\infty$};
			      \node[above] at (h.north) {$\infty$};
			      \node[below] at (g.south) {$-2$};
			      \node[below] at (f.south) {$-1$};

			      \draw[->, bend left=15, >=stealth, line width=1pt] (e) to node[above] {-1} (b);
			      \draw[->, bend left=15, >=stealth, line width=1pt] (b) to node[above] {4} (a);
			      \draw[->, >=stealth, line width=1pt] (e) to node[above] {-2} (g);
			      \draw[->, >=stealth, line width=1pt] (h) to node[below] {-3} (g);
			      \draw[->, >=stealth, line width=1pt] (g) to node[below] {1} (f);
			      \draw[->, bend left=15, >=stealth, line width=1pt] (f) to node[left] {2} (a);
			      \draw[->, bend right=15, >=stealth, line width=1pt] (d) to node[left] {4} (f);
			      \draw[->, bend right=15, >=stealth, line width=1pt] (f) to node[right] {-3} (d);
		      \end{tikzpicture}
	      \end{center}

	      Шаг 3:
	      \begin{center}
		      \begin{tikzpicture}
			      \node[draw, circle] (e) at (0,2) {e};
			      \node[draw, circle] (b) at (2,2) {b};
			      \node[draw, circle] (a) at (4,2) {a};
			      \node[draw, circle] (d) at (6,2) {d};
			      \node[draw, circle] (h) at (0,0) {h};
			      \node[draw, circle] (g) at (3,0) {g};
			      \node[draw, circle] (f) at (6,0) {f};

			      \node[above] at (e.north) {$0$};
			      \node[above] at (b.north) {$-1$};
			      \node[above] at (a.north) {$1$};
			      \node[above] at (d.north) {$-4$};
			      \node[above] at (h.north) {$\infty$};
			      \node[below] at (g.south) {$-2$};
			      \node[below] at (f.south) {$-1$};

			      \draw[->, bend left=15, >=stealth, line width=1pt] (e) to node[above] {-1} (b);
			      \draw[->, bend left=15, >=stealth, line width=1pt] (b) to node[above] {4} (a);
			      \draw[->, >=stealth, line width=1pt] (e) to node[above] {-2} (g);
			      \draw[->, >=stealth, line width=1pt] (h) to node[below] {-3} (g);
			      \draw[->, >=stealth, line width=1pt] (g) to node[below] {1} (f);
			      \draw[->, bend left=15, >=stealth, line width=1pt] (f) to node[left] {2} (a);
			      \draw[->, bend right=15, >=stealth, line width=1pt] (d) to node[left] {4} (f);
			      \draw[->, bend right=15, >=stealth, line width=1pt] (f) to node[right] {-3} (d);
		      \end{tikzpicture}
	      \end{center}

	      После 4-го шага значения в вершинах графа
	      изменяться не будут

	\item Можно использовать алгоритм Беллмана-Форда,
	      но так как про граф известно, что все рёберные длины
	      не превышают $k$, то всего итраций понадобится $k + 1$
	      (ещё одна итерация для проверки наличия циклов отрицательного).

	      Алгоритм:
	      \begin{lstlisting}
def update_dist(
    graph: dict[object, list[tuple[float, object]]],
    dist: dict[object, float],
):
    for node, neighbors in graph.items():
        for weight, neighbor in neighbors:
            if dist[node] + weight < dist[neighbor]:
                dist[neighbor] = dist[node] + weight
                return True
    return False


def k_bellman_ford(
    graph: dict[object, list[tuple[float, object]]],
    source: object,
    k: int,
) -> dict[object, float]:
    # assume graph is strongly connected
    dist = {node: float("inf") for node in graph}
    dist[source] = 0
    for _ in range(k):
        changed = update_dist(graph, dist)

        if not changed:
            return dist

    changed = update_dist(graph, dist)

    if changed:
        raise ValueError("Negative cycle detected")

    return dist
          \end{lstlisting}

	      Сложность алгоритма $O(k |E|)$, так как
	      на каждой из $k$ итераций происходит обновление
	      всех $|E|$ рёбер.

	\item Для любой вершины дерева её независимым множеством
	      будут те вершины, которые расположены на уровне
	      такой же чётности, как и эта вершина (например, для корня на
	      уровне 0 - все вершины на четных уровнях). Это можно доказать
	      рекурсивным раскрашиванием вершин дерева в два цвета:
	      независимо от того, с какой вершины начать, можно покрасить
	      всех соседей каждой вершины в один цвет, отличный от цвета
	      самой вершины, в противном случае в дереве есть цикл.

	      Тогда нам нужно рассмотреть только два множества вершин:
	      на четных уровнях и на нечетных.

	      Алгоритм:
	      \begin{lstlisting}
from collections import deque


def find_largest_independent_set(tree: dict[object, list[object]]) -> set[object]:
    even = set()
    odd = set()

    start_node = next(iter(tree))
    even.add(start_node)

    queue = deque([(start_node, None, 0)])
    while len(queue) > 0:
        node, parent, level = queue.popleft()

        if level == 0:
            even.add(node)
        else:
            odd.add(node)

        for child in tree[node]:
            if child == parent:
                continue

            queue.append((child, node, (level + 1) % 2))

            if child in even or child in odd:
                raise ValueError("Graph is not a tree")

    if len(even) > len(odd):
        return even
    return odd
    \end{lstlisting}
	      Сложность - $O(|V| + |E|) = O(|V|)$, так как алгоритм обходит дерево,
	      посещая каждую вершину один раз.

	\item За $O(|V| + |E|)$ задачу можно свести к обычному
	      поиску в ширину, совмещая вершины, соединённые
	      ребром с весом $0$, далее за $O(|V| + |E|)$ можно
	      найти кратчайший путь от $s$ до каждой вершины
	      алгоритмом \texttt{BFS}.

	      Алгоритм:
	      \begin{lstlisting}
from collections import deque


def get_graph_groups(
	graph: dict[object, list[tuple[int, object]]],
) -> dict[object, set[object]]:
	united_nodes = {}

	for node, neighbors in graph.items():
		for weight, neighbor in neighbors:
			if weight == 0:
				united_nodes[node] = neighbor
				united_nodes[neighbor] = node

	grouped_nodes = {}
	for node, neighbor in united_nodes.items():
		grouped_nodes[node] = grouped_nodes.get(node, set([node])).union([neighbor])

	return grouped_nodes


def shortest_distances(
	graph: dict[object, list[tuple[int, object]]],
	source: object,
) -> dict[object, int]:
	groups = get_graph_groups(graph)

	distances = {}
	queue = deque([(source, 0)])

	while queue:
		node, cur_len = queue.popleft()

		if node in distances:
			continue

		for subnode in groups.get(node, set([node])):
			distances[subnode] = cur_len

		for weight, neighbor in graph[node]:
			if weight == 0:
				continue
			if weight != 1:
				raise ValueError("Graph does not satisfy the requirements")

			if neighbor in distances:
				continue

			queue.append((neighbor, cur_len + 1))

	return distances
	\end{lstlisting}
	      Сложность - $O(|V| + |E|)$, так как сначала алгоритм обходит
	      все вершины и рёбра графов за $O(|V| + |E|)$ для поиска
	      групп вершин, соединённых рёбрами с весом $0$, а затем
	      запускает обход в ширину для поиска кратчайших путей от
	      данной вершины до всех остальных.

\end{enumerate}

\end{document}


